\section{Agent-Analyse (Kap. 01)}
\label{sec:agent-analyse}

\subsection{PEAS-Analyse}

Ein Agent ist durch seine Wahrnehmung und Wirkung in einer Umgebung bestimmt,
PEAS zwingt dabei zur expliziten Operationalisierung des Erfolgs
\parencite[Kap.~2]{russellNorvig2021}. Für den Tutor-Workflow ergibt sich
\cref{tab:peas}.

\begin{table}[htbp]
  \centering
  \small
  \begin{tabularx}{\linewidth}{@{}lX@{}}
    \toprule
    PEAS-Komponente & Projektspezifische Ausprägung und Codebeleg \\
    \midrule
    Performance & Fachlich plausible, nicht leere und zur Hilfestufe passende
    Antwort, Erfüllung von \texttt{success\_criteria}, korrekte Werkzeugwahl
    und Reviewer-Freigabe. Eine einzige numerische Nutzenfunktion existiert
    nicht. Die separaten Evaluationsmetriken prüfen sokratische Führung,
    Vermeidung direkter Lösungen, Trace-ID und Tool-Abdeckung. \\
    Environment & Authentifizierte Lernende, Angular-Oberfläche, aktuelle
    Aufgaben und Dokumente, PostgreSQL-Zustand, LLM-Anbieter sowie Resultate
    interner Werkzeuge. \\
    Actuators & Tutor-Text als primäre Wirkung, zusätzlich die im Chat
    tatsächlich freigebbaren Werkzeuge \texttt{evaluate\_answer},
    \texttt{award\_coins\_and\_exp} und \texttt{add\_calendar\_note}. Die
    letzten beiden verändern persistenten Zustand. \\
    Sensors & Nachricht, Zeit, Dialoghistorie, aktive Aufgabe, Hilfestufe,
    Frontend-Kontext, bis zu drei semantisch ähnliche Dokumentchunks und
    Werkzeugergebnisse. Die Bildung des \texttt{AgentContext} erfolgt im
    Endpunkt \path{backend/app/api/v1/endpoints/agent.py}. \\
    \bottomrule
  \end{tabularx}
  \caption{PEAS-Analyse des agentischen Tutor-Workflows.}
  \label{tab:peas}
\end{table}

\subsection{Taxonomie und Umgebungsdimensionen}

StudyMummy ist am treffendsten ein \emph{begrenzt modellbasierter,
zielorientierter und deliberativer} Softwareagent. Modellbasiert ist er, weil
Sitzung, Aufgabe, Hilfestufe und Historie einen internen Zustand bilden.
Zielorientiert ist er durch das je Turn erzeugte \texttt{objective} und ein bis
drei Erfolgskriterien. Der JSON-Plan und die nachfolgende Ausführung gehen über
einen Reflexagenten hinaus. Eine Einordnung als Utility-Agent wäre dagegen zu
stark: Weder werden konkurrierende Pläne bewertet noch wird erwarteter Nutzen
maximiert. Auch ein Lernagent im klassischen Sinn ist nicht vollständig
implementiert. Zwar existieren Lernprofil und Update-Werkzeug, doch die
Capability-Policy gibt dieses Werkzeug im Chat nie frei und das Profil fließt
nicht in den \texttt{AgentContext} ein.

Die Umgebung ist nur \emph{partiell beobachtbar}, da Wissen, Motivation und
tatsächliches Verständnis der lernenden Person indirekt aus Text erschlossen
werden. Sie ist \emph{sequenziell} und \emph{dynamisch}: Frühere Antworten,
Aufgabenstatus und Dokumente beeinflussen spätere Entscheidungen und können
sich zwischen Turns ändern. Datenbankoperationen sind überwiegend
deterministisch, LLM-Ausgaben und Nutzerverhalten jedoch stochastisch. Aktionen
und Zustände sind in der Implementierung diskret typisiert, der relevante
Zustandsraum bleibt wegen freier Sprache sehr groß. Die eigenen
Werkzeugtransitionen sind bekannt, die Reaktion von Mensch und Modell nur
teilweise. Innerhalb des Systems interagieren mehrere logische Agentenrollen,
gegenüber der Lernumgebung ist der Workflow zugleich ein gemeinsamer
Tutorakteur.
